\documentclass[12pt,a4paper]{article}

\usepackage[margin=1in]{geometry}
\usepackage{booktabs}
\usepackage{longtable}
\usepackage{array}
\usepackage{graphicx}
\usepackage{hyperref}
\usepackage{xcolor}
\usepackage{enumitem}

\title{Design Pattern Analysis in a FastAPI-Based Logistics Management Application}
\author{Design Pattern Case Study}
\date{October 2026}

\begin{document}

\maketitle

\tableofcontents
\newpage

\section{Application Overview}

This case study analyzes the Warehouse \& Transportation Management System (WTMS), a FastAPI-based logistics management application using PostgreSQL.

The application provides warehouse management, multi-warehouse inventory, storage zones and bins, stock movements, orders, vehicles, drivers, routes, shipments, shipment tracking, and transportation management.

The baseline application uses a repository-based data-access architecture.

\section{Environment}

The analyzed environment was verified as follows:

\begin{itemize}
    \item Backend: FastAPI / Python
    \item Python: 3.11.5
    \item Database: PostgreSQL 18.6
    \item Java: OpenJDK/Temurin 21.0.1 LTS
    \item JavaScript runtime: Node.js 22.15.1
    \item C++: MinGW.org GCC 6.3.0 with C++17 and \texttt{-O2}
\end{itemize}

\section{Part 1: Baseline Pattern Analysis}

\subsection{Confirmed Repository Pattern}

The primary confirmed design pattern is the Repository Pattern.

Four repository instances were confirmed from source-code evidence:

\begin{enumerate}
    \item WarehouseRepository
    \item InventoryRepository
    \item OrderRepository
    \item TransportationRepository
\end{enumerate}

The evidence shows API endpoints obtaining repository instances through FastAPI dependency injection and repositories encapsulating database operations through the Database abstraction.

FastAPI dependency injection and the Database abstraction were treated as supporting mechanisms rather than counted as additional design-pattern instances.

No additional pattern was counted without sufficient source-code evidence.

\subsection{Pattern Evidence}

The detailed source evidence is maintained in:

\begin{itemize}
    \item \texttt{CaseStudy/Part1/application\_analysis.txt}
    \item \texttt{CaseStudy/Part1/repository\_pattern\_evidence.txt}
    \item \texttt{CaseStudy/data/patterns.csv}
    \item \texttt{CaseStudy/data/pattern\_counts.csv}
\end{itemize}

\section{Part 2: Design Evolution}

\subsection{Multiple Warehouses}

Multiple-warehouse support already exists in the baseline application.

The baseline database initialization was verified with four warehouse records. The \texttt{warehouses} table and warehouse references in inventory-related entities provide database-level support for multiple warehouses.

No artificial change was introduced for this requirement. It is therefore documented as a baseline capability.

Evidence is maintained in:

\texttt{CaseStudy/Part2/MultiWarehouse\_Baseline/multi\_warehouse\_baseline.txt}

\subsection{Change 1: Delivery Partner}

A new \texttt{delivery\_partners} table was introduced.

The \texttt{shipments} table was extended with \texttt{delivery\_partner\_id}, together with an index and foreign-key relationship.

The application was extended with delivery-partner models, repository operations, and API endpoints.

Runtime verification successfully created and retrieved a delivery partner and created a shipment associated with the partner.

Detailed evidence is maintained in:

\begin{itemize}
    \item \texttt{CaseStudy/Part2/Change1\_DeliveryPartner/change1\_analysis.txt}
    \item \texttt{CaseStudy/Part2/Change1\_DeliveryPartner/change1\_schema.sql}
    \item \texttt{CaseStudy/Part2/Change1\_DeliveryPartner/change1\_delivery\_partner.puml}
\end{itemize}

\subsection{Change 2: Real-Time Tracking}

A new \texttt{shipment\_tracking} table was introduced.

The table stores shipment ID, latitude, longitude, location, status, and recording timestamp.

Tracking creation and tracking-history retrieval were implemented through the transportation repository and API.

Runtime verification successfully inserted tracking updates and retrieved shipment tracking history.

Detailed evidence is maintained in:

\begin{itemize}
    \item \texttt{CaseStudy/Part2/Change2\_Tracking/change2\_analysis.txt}
    \item \texttt{CaseStudy/Part2/Change2\_Tracking/change2\_schema.sql}
    \item \texttt{CaseStudy/Part2/Change2\_Tracking/change2\_source\_evidence.txt}
    \item \texttt{CaseStudy/Part2/Change2\_Tracking/change2\_tracking.puml}
\end{itemize}

\subsection{Change 3: Route Changes}

A new \texttt{shipment\_route\_history} table was introduced to preserve previous and new route assignments.

The repository operation performs the route change transactionally:

\begin{enumerate}
    \item Lock and read the current shipment route.
    \item Update the shipment route.
    \item Insert the previous and new route IDs into route history.
\end{enumerate}

Runtime verification changed shipment 4 from route 1 to route 2 and recorded the corresponding route-history row.

Detailed evidence is maintained in:

\begin{itemize}
    \item \texttt{CaseStudy/Part2/Change3\_RouteChanges/change3\_analysis.txt}
    \item \texttt{CaseStudy/Part2/Change3\_RouteChanges/change3\_schema.sql}
    \item \texttt{CaseStudy/Part2/Change3\_RouteChanges/change3\_source\_evidence.txt}
    \item \texttt{CaseStudy/Part2/Change3\_RouteChanges/change3\_route\_changes.puml}
\end{itemize}

\subsection{Pattern Evolution}

The confirmed Repository Pattern instances were tracked across the baseline and three changes.

The four baseline repository instances remain present. TransportationRepository is marked as extended across the changes because the transportation-related functionality was expanded.

The complete evolution record is stored in:

\texttt{CaseStudy/data/patterns\_evolution.csv}

\section{Part 3: Cross-Language Implementation}

The Repository Pattern was implemented independently in Java, Python, JavaScript, and C++.

Each implementation performs equivalent shipment operations:

\begin{itemize}
    \item Save a shipment
    \item Find a shipment by ID
    \item Retrieve all shipments
    \item Verify the expected result
\end{itemize}

All four implementations were successfully executed.

\subsection{Measured Results}

\begin{center}
\begin{tabular}{lrrrr}
\toprule
Language & Lines & Complexity & Median (ms) & Result \\
\midrule
Java & 74 & 3 & 81.214 & PASS \\
Python & 45 & 1 & 93.486 & PASS \\
JavaScript & 49 & 3 & 46.009 & PASS \\
C++ & 65 & 3 & 14.209 & PASS \\
\bottomrule
\end{tabular}
\end{center}

The detailed measurements, including all five runs for every language, are stored in:

\texttt{CaseStudy/data/cross\_language.csv}

\subsection{Timing Methodology}

For every implementation:

\begin{enumerate}
    \item The implementation was compiled where compilation was required.
    \item One warm-up execution was performed.
    \item Five measured executions were performed.
    \item PowerShell \texttt{Measure-Command} was used.
    \item The median of the five measured executions was recorded.
\end{enumerate}

The measurements are process-level timings and include process startup overhead. They should not be interpreted as general programming-language performance rankings.

\subsection{Complexity}

Cyclomatic complexity was documented using the McCabe definition:

\[
CC = 1 + \text{number of decision points}
\]

The values were based on the actual control-flow constructs in the implementations.

\section{Evidence and Data Management}

Pattern instances were counted only when source-code evidence supported the pattern.

Class names, annotations, README descriptions, and LLM suggestions were not treated as sufficient evidence by themselves.

The case study data files are:

\begin{itemize}
    \item \texttt{patterns.csv}
    \item \texttt{pattern\_counts.csv}
    \item \texttt{patterns\_evolution.csv}
    \item \texttt{cross\_language.csv}
    \item \texttt{llm.csv}
\end{itemize}

The \texttt{llm.csv} file contains only its header at the time of report generation because fabricated prompts or outputs were not added.

\section{Limitations}

The cross-language programs are small representative implementations of the Repository Pattern and are not replacements for the complete WTMS backend.

Process-level timing includes runtime startup overhead.

The measurements are specific to the local environment and these demonstration programs.

The case study identifies patterns from available source-code evidence and does not infer patterns solely from naming or documentation.

\section{Conclusion}

The baseline WTMS application demonstrates a Repository-based data-access architecture with four confirmed repository instances.

The evolution scenarios demonstrate how the existing transportation repository was extended to support delivery partners, shipment tracking, and route-history management while preserving the repository abstraction.

The cross-language implementations demonstrate the same Repository Pattern structure in Java, Python, JavaScript, and C++, with reproducible functional tests and measured process-level execution times.

\end{document}

